\documentclass{jsarticle}
\usepackage[dvipdfm,final]{graphicx}
\usepackage{chemfig}
\usepackage{amsmath}
\usepackage{amssymb}
\begin{document}
\title{Recycleにおいての塩化ビニルの \\
ナトリウムプロパノキシドでの分解}
\author{睦世姉さんの論文について}
\date{}
\maketitle
  $$ \chemfig{\mathrm{[}-CH_2 -C=N-CH_2-\mathrm{]}-Cl} + 
  \chemfig{*6(-(-[180]NaO)-=-=-=)} + \text{熱} 
   \to \chemfig{CO_2 + H_2O + NaCl} $$
   $$ \chemfig{\mathrm{[}-CH_2-C=N-CH_2-\mathrm{]}-Cl} +
   \chemfig{NaO-CH_2-C([:90]=O)-CH_2-CH_3} + \text{熱}
   \to \chemfig{CO_2 + H_2O + NaCl} $$
   この化学式から推算式でもある、プラスチックと二酸化ケイ素、塩化ビニール
   が分解されるのが、真逆の人が服用すると、プラバスタチンとスリムドガン
   と同じく、推算式の数値が上がる。僧侶と同じ瞑想法と原理が
   一致している。
   $$ \chemfig{*6(-=-(-COOH)=-=-C=C-C(-C=C-C)=C-C=C(-H)-OH)} $$
   $$ \chemfig{*5(-=-C-N=C-OH)} $$
   $$ \chemfig{*6(=-(-[,1.2]C?(=[:-90,1.2]O))=
   (-[,1.2]C(=[:90,1.2]O)-[:-40,1.6]O?)-=-)} $$
   これらの構造式は、それぞれ、ヒドロキシ基、カルボニル基、カルボキシル基、
   アミン基、ギ酸、アセチル基、アセト基、グリシン、アミノ酸の２０種類の
   一基、ビタミンCはシス・トランス型があり、松果体に作用して、
   体内の血流を調節する。$ \chemfig{-C=N-Cl} $の真逆の作用をする。
   アデノシンは、同じく血流を調節して、エピネフリンはエンドルフィン
   と同じく痛み止めに効く。ノンアドレナリンは、無水フタル酸と同じ作用をする。
   グリシンには、２０種類のアミノ酸の中で唯一L,R基が無い。不斉炭素原子が
   無い。
   推算式の元を減らすか増やすのに、血圧と血流、体内の精神作用のドーパミンの
   調節をeGFRの機能を上げると、結果、すべての精神作用物質が調節できる。:





\end{document}
